\scanpage{148}\mspage{3}
for having corrected in my mind a gross error, due to my great ignorance of algebraic topology and of homotopy -- I was indeed under the impression that $H$-spaces satisfying fairly strict conditions of associativity and commutativity (say equivalent to $\unsure{{}'}\!\Omega^{i}X$ with $i$ arbitrarily large) corresponded to\unsure{.} \unsure{commutative|commutative} topological groups (drawing on a few analogies \dots). So entirely in agreement with your observations page~5.\unsure{.}

On the other hand, I am still intrigued by the following question: is\unsure{.} \unsure{there|there} not an analogue of the Dold-Puppe carpet for ($1/2$) simplicial groups (not necessarily commutative) and what form does it take? To tell the truth, I regard the yoga
\diagram{\begin{tikzcd}[ampersand replacement=\&]
\text{set }1/2\text{ simplicial}\arrow[leftrightarrow,r] \& \infty\text{-}\deleted{\unsure{\text{category}}}\text{groupoid}
\end{tikzcd}}
as being essentially the ultimate ``set-theoretic'' version of Dold-Puppe, which would follow from it by merely making explicit that the abelian groups in $(\infty\text{-cat})$ ``are none other'' than the chain complexes in $(\mathrm{Ab})$. It would therefore be necessary to determine what the groups of $(\infty\text{-cat})$ are: I can tell you how matters stand in $(1\text{-cat})$ -- this will be \unsure{develop|divulg}%
% The manuscript has a terminal line-end hyphen after the incomplete final word.
% archive sheet number: 148
